This survey examines the two Kakeya manuscripts released by OpenAI on September 23 and 24, 2026. The first claims to prove the three-dimensional Kakeya maximal conjecture, sharpening the density exponent $\lambda^{K(\varepsilon)}$ left open after Wang--Zahl to the critical $\lambda^{3}$; the second claims that every four-dimensional Kakeya set has Hausdorff dimension four, with no measurability, compactness, or stickiness hypothesis. Neither manuscript has been peer-reviewed or formalized in Lean, and the repository states that unformalized results may contain issues. The survey proceeds in four parts. First, a verification file: timeline, manuscript parameters, dependency chain, and verification status, including the observation that the four-dimensional result imports a weighted plank bound from its three-dimensional companion, so its validity is coupled to that of the companion. Second, a technical reconstruction: the reduction from an arbitrary Kakeya set to a weighted model, the bookkeeping of a positive ``horizon'' (a time-cost obstruction), quadratic charts with mass constraints, extremal critical paths, and the three narrowness regimes. Third, a provenance analysis that classifies each core technique as direct inheritance, in-school modification, cross-discipline import, or first systematic use within the Kakeya lineage; comparison with Wang--Zahl and Guth--Wang--Zahl shows that the entropy ledger, probabilistic coupling, quantitative real algebraic geometry, and Fourier-decay estimates have no precedent in the traditional Kakeya toolkit. Fourth, a reproducible verification roadmap and directions for future work along four tracks: verification, mathematics, methodology transfer, and research workflow. Claims of the manuscripts are kept strictly separate from established theorems, and each judgement carries a stated confidence level.